%% kapa: exam=1 round=32 year=2021 period=1 form=A subject=civil_law subject_ko=민법 seq=37 number=37
%% source: Q-Net 감정평가사 기출문제 (cst00309) artlSeq=5208397 file=extracted/1008/5208397/A형/2021년도 제32회 감정평가사 1차 1교시 A형.pdf
%% conversion: kordoc 4.24.51 -> tools/gen1.py (auto); review=proofread-sonnet; flags=none
\begin{question}{37}
질권에 관한 설명으로 옳지 않은 것은? (다툼이 있으면 판례에 따름)
\begin{choices}
  \item 양도할 수 없는 물건은 질권의 목적이 되지 못한다.
  \item 질권자는 채권의 변제를 받기 위하여 질물을 경매할 수 있다.
  \item 채권질권의 효력은 질권의 목적이 된 채권 외에 그 채권의 지연손해금에는 미치지 않는다.
  \item 질권의 목적인 채권의 양도행위는 질권자의 이익을 해하는 변경에 해당되지 않으므로 질권자의 동의를 요하지 않는다.
  \item 수개의 채권을 담보하기 위하여 동일한 동산에 수개의 질권을 설정한 경우, 그 순위는 설정의 선후에 의한다.
\end{choices}
\end{question}
\answer{3}
